\subsection*{Limitation}
While the EDCKD model shows promise, several limitations must be acknowledged regarding its deployment. A key concern is the small sample size of the development cohort (n=491, with only 56 CKD events), which limits the statistical power for detecting rare subgroup effects and may affect model stability. The limited number of CKD cases resulted in wide confidence intervals for performance metrics, particularly precision estimates. To address this, future studies should incorporate larger multicenter cohorts with diverse patient populations during training. The model's training primarily on the UAE Clinical CKD dataset may limit its generalizability to diverse chronic kidney disease populations encountered in different geographic regions and healthcare settings, as evidenced by the performance drop from 96.97\% to 80.50\% accuracy on external validation. The model should be validated on a wider array of datasets that reflect various demographics, laboratory protocols, and clinical practices. Despite attempts to mitigate class imbalance through BorderlineSMOTE, this issue remains, particularly given the 11.4\% CKD prevalence in the development dataset. This imbalance can significantly affect performance metrics such as recall and precision. For instance, in Table 6, the proposed Random Forest model demonstrates perfect precision (100\%) on the test set, but this is based on only 8 correct CKD predictions, resulting in wide confidence intervals. Similarly, the external validation showed a disparity between the internal and external performance metrics, indicating that the model may excel in the development population while underperforming in independent cohorts. Future research should explore strategies such as cost-sensitive learning, advanced resampling techniques, and ensemble methods to better address this imbalance and improve the overall performance of the model across diverse clinical settings. The model relies on feature availability and measurement standardization, which posed challenges during external validation, where only 9 of the 23 selected features were available in the UCI dataset. Missing or incompatible features significantly affect prediction reliability, posing challenges for cross-institutional deployment. Ensuring feature completeness and developing robust imputation strategies are crucial for future iterations of the EDCKD system. Additionally, the temporal gap in the training data (collected 2008--2017) may limit its applicability to current clinical practices, as diagnostic protocols, treatment guidelines, and patient demographics have evolved. Prospective validation using contemporary data is required to confirm sustained performance. By addressing these limitations, we can improve the robustness and applicability of the EDCKD model in real-world clinical settings.
